
\section{Related Work}
\label{section: related work}
In this section, we present the existing related work in the field of blockchain-enabled IoT. Then, we review research on blockchain sharding techniques, including random-based sharding, community-based sharding and trust-based sharding. Additionally, we compare the use of DRL technology in sharded blockchains.

\subsection{Blockchain and IoT}
With the rapid development of blockchain technology, blockchain-enabled IoT has become more secure and reliable. Existing blockchain-enabled IoT systems primarily focus on framework design and consensus protocol design. For instance, Kang \textit{et al.}~\cite{kang2018blockchain} proposed a solution to address security and privacy challenges in Vehicular Edge Computing and Networks (VECONs) based on consortium blockchain and smart contracts. In the medical scenario, Gadekallu \textit{et al.}~\cite{gadekallu2021blockchain} presented a blockchain-based solution to enhance the security of datasets generated from IoT devices in e-health applications. The proposed solution utilizes a blockchain platform and a private cloud to secure and verify the datasets, reducing tampering risks and ensuring reliable results. 
Xu \textit{et al.} proposed a blockchain-enabled system for data provenance with 5G and LoRa network~\cite{9526860,YU2021102492}.


\subsection{Blockchain Sharding}

\subsubsection{Random-based sharding}
In the realm of blockchain sharding research, a series of techniques based on random sharding have been proposed to enhance system scalability and fault tolerance. Elastico~\cite{luu2016secure} is a pioneering protocol that introduced sharding within a permissionless setting, dynamically scaling the network in tandem with the increasing number of participating nodes. OmniLedger ~\cite{kokoris2018omniledger} took this further, ensuring linear scalability through an effective sharding mechanism that seamlessly manages cross-shard transactions. RapidChain ~\cite{zamani2018rapidchain} built upon these concepts with a full sharding solution that streamlines the partitioning of blockchain states and the processing of transactions. Monoxide~\cite{wang2019monoxide} introduced asynchronous consensus areas to boost transaction efficiency, though its static sharding strategy may fall short under dynamically changing network conditions. These innovations have significantly contributed to system robustness by improving randomness and fault tolerance, yet challenges such as frequent node reassignments remain, leading to non-trivial system overheads.

\subsubsection{Community-based sharding}
Community-based sharding is a partitioning method in blockchain networks that divides the system into smaller subsets or shards based on the interactions between nodes. Nodes and transactions are depicted as a graph and are partitioned into smaller subgraphs or shards, each comprising a subset of nodes and transactions. Fynn \textit{et al.}~\cite{fynn2018challenges} investigate the Ethereum blockchain's scalability through sharding implementation. They model the Ethereum blockchain as a graph and analyze different algorithms for graph partitioning, such as the Kernighan-Lin algorithm ~\cite{kernighan1970efficient} and METIS~\cite{karypis1998fast}. Zhang \textit{et al.}~\cite{zhang2022community} proposed community detection-based sharding to enhance sharded blockchain scalability. They used the community detection algorithm to group nodes with frequent trading in the same shard to reduce CSTs. However, this solution demands significant computational resources and high communication costs.


\subsubsection{Trust-based sharding}
Trust-based blockchain sharding divides the network into smaller shards based on node trust. Nodes are evaluated by reputation and behavior, grouping trusted nodes to enhance security and performance and evenly distributing dishonest nodes. Yun \textit{et al.}~\cite{yun2019trust} presented a Trust Based Shard Distribution (TBSD) scheme to enhance system security and prevent collusion, specifically aiming to address the $1\%$ attack. However, their approach does not consider shard load balance and trust difference, potentially leading to system delays. Huang \textit{et al.}~\cite{huang2020repchain} proposed RepChain, a reputation-based blockchain system with sharding for high throughput, security, and node cooperation. It utilizes a double-chain architecture: transaction chain and reputation chain. However, the double-chain architecture chain increases system complexity and resource needs, resulting in additional overhead. Zhang \textit{et al.}~\cite{zhang2023optimized} proposed a new blockchain sharding model to enhance security and efficiency. Their approach considers shard trust, latency, and node count differences, reducing the risk of blockchain failure. However, the effectiveness of their proposed model heavily depends on obtaining accurate information, which can be challenging to obtain in dynamic blockchain sharding scenarios.
 

\subsection{Deep Reinforcement Learning-based Sharding}
DRL-based sharding solutions have been making strides in the realm of IoT. Liu \textit{et al.}~\cite{liu2019performance} were pioneers, leveraging DRL in a blockchain framework tailored for Industrial IoT (IIoT). While they dynamically adjusted critical parameters, they missed addressing dishonest attacks. Similarly, another work by Liu \textit{et al.}~\cite{liu2018blockchain} harnessed Ethereum and DRL for IIoT data security but sidestepped throughput scalability concerns. On the other hand, Qiu \textit{et al.} presented service-oriented blockchain solutions, using DRL for refined block production and bandwidth allocation~\cite{qiu2019service,qiu2018blockchain}. However, these lacked centralized security measures, making them susceptible to dishonest threats.

Further innovations came from Yun \textit{et al.}~\cite{yun2020dqn}, who introduced a DQN-optimized framework for sharded blockchains. Despite its advancements, the approach overlooked intricate attack strategies in DRL-based sharding and was confined by its central Q-learning reliance. Meanwhile, Yang \textit{et al.}~\cite{yang2022sharded} incorporated K-means clustering in a sharded blockchain, utilizing DRL for optimization. Yet, their methodology was marred by the computational intensity and intricacy of DRL.
Most existing articles~\cite{liu2019performance, liu2018blockchain, qiu2019service, qiu2018blockchain} lack effective analysis of dishonest attacks in IoT using DRL-based blockchain, neglecting the complexities and depth of practical security threats. In contrast, our proposed model analyses the $1\%$ attack problem in the blockchain sharding model and considers strategic collusion attacks involved in existing related studies. By understanding attackers' motivation, we develop proactive defense mechanisms and risk mitigation strategies. In this paper, we design a trustworthy DRL-based blockchain sharding system under IoT scenarios. 




